\chapter{杂项}

\section{自定义运算符}

\subsection{结构体内重载比较、加法与输出}

记录按分数降序、罚时升序确定优先级。

\begin{itemize}
    \item 二元成员运算符的左操作数为 \texttt{*this}。以 \texttt{const} 限定成员函数后可比较常量对象。
    \item \texttt{+=} 修改左操作数并返回引用。\texttt{+} 复制左操作数后复用 \texttt{+=}。
    \item 在结构体内用 \texttt{friend} 定义流输出。以流为左操作数，返回流引用以支持连续输出。
\end{itemize}

\cppfile{code/zx/operator_overload.hpp}

默认排序将高优先级记录置于末尾，默认优先队列将其置于堆顶。\texttt{greater<Score>} 反转次序。

比较运算符须满足严格弱序。累加结果须能用 \texttt{ll} 表示。

\subsection{128 位整数流输出}

让代码中的 \texttt{i128} 对应有符号 128 位类型后，可用 \texttt{cout << x} 输出十进制值。\texttt{x} 取最小负值时改用第八章的 \texttt{toString(x)}。

\cppfile{code/zx/zx1.cpp}

\section{对拍}

\subsection{线下赛运行脚本}

执行 \texttt{sh run.sh solution sample}，传入程序名和样例名。脚本编译程序，读取样例输入并与对应的期望输出比较。

\cppfile[language=bash,numbers=none]{code/zx/run.sh}

\subsection{循环对拍脚本}

准备名为 \texttt{gen.py} 的生成器和两个可执行程序 \texttt{std}、\texttt{sol}，分别充当参考解和待测解。执行 \texttt{sh stress.sh}，输出不同时结束，查看保留的输入和两份输出。

\cppfile[language=bash,numbers=none]{code/zx/stress.sh}

\section{随机数与随机化}

\subsection{随机化}

调用 \texttt{rd(l,r)} 从闭区间 $[l,r]$ 取随机整数。

\cppfile{code/zx/sjh.cpp}

\subsection{模拟退火}

构造 \texttt{SA sa(n)} 并填入非空点集的 \texttt{xs} 和 \texttt{ys}。

调用 \texttt{run(init\_T,end\_T,alpha,iter\_per\_T)}，从 \texttt{best\_x}、\texttt{best\_val} 读取位置与目标值。用于其他题目时按目标改写 \texttt{energy}。

\cppfile{code/zx/SA.cpp}
